\documentclass[preview]{standalone}

\usepackage{amsmath}
\usepackage{amssymb}
\usepackage{stellar}
\usepackage{definitions}

\begin{document}

\id{primitive-polynomials}
\genpage

\section{Primitive polynomials}

Let \(A\) be a \uniquefactorizationdomain and let
\[
    K=\quotientfield(A).
\]
We identify \(A[x]\) with its canonical subring in \(K[x]\).

\begin{snippetdefinition}{primitive-polynomial-definition}{Primitive polynomial}
    Let \(A\) be a \uniquefactorizationdomain. A nonzero \polynomial
    \(f\in A[x]\) is \emph{primitive} if
    \[
        \nu_p(f)=0
        \qquad
        \text{for every \primeelement \(p\in A\)}.
    \]
    Equivalently, no prime element of \(A\) divides every coefficient of
    \(f\).
\end{snippetdefinition}

\begin{snippetproposition}{monic-polynomials-primitive}{Monic polynomials are primitive}
    Let \(A\) be a \uniquefactorizationdomain. Every monic \polynomial
    in \(A[x]\) is primitive.
\end{snippetproposition}

\begin{snippetproof}{monic-polynomials-primitive-proof}{monic-polynomials-primitive}{Monic polynomials are primitive}
    If
    \[
        f=a_0+a_1x+\cdots+a_nx^n
    \]
    is monic, then \(a_n=1_A\). Hence
    \(\nu_p(a_n)=\nu_p(1_A)=0\) for every \primeelement \(p\), and
    therefore \(\nu_p(f)=0\).
\end{snippetproof}

\begin{snippetexample}{primitive-integer-polynomial-example}{A primitive integer polynomial}
    The \polynomial
    \[
        f=2x^3+3x^2+5x\in\integers[x]
    \]
    is primitive because no prime integer divides all of its
    coefficients. In valuation language,
    \[
        \nu_p(f)=0
        \qquad
        \text{for every prime integer }p.
    \]
\end{snippetexample}

\begin{snippetexample}{rational-polynomial-factorization-example}{A factorization over the rationals}
    In \(\rationalnumbers[x]\),
    \[
        (x^2+1)^2(x-3)
    \]
    is a product of irreducible polynomials: \(x-3\) is linear, and
    \(x^2+1\) has no root in \(\rationalnumbers\).
\end{snippetexample}

\section{Gauss's lemma}

\begin{snippetlemma}{primitive-polynomial-product-lemma}{Products of primitive polynomials}
    Let \(A\) be a \uniquefactorizationdomain. The product of two
    primitive polynomials in \(A[x]\) is primitive.
\end{snippetlemma}

\begin{snippetproof}{primitive-polynomial-product-lemma-proof}{primitive-polynomial-product-lemma}{Products of primitive polynomials}
    If \(f,g\in A[x]\) are primitive, then for every \primeelement
    \(p\in A\),
    \[
        \nu_p(fg)=\nu_p(f)+\nu_p(g)=0.
    \]
    Thus \(fg\) is primitive.
\end{snippetproof}

\begin{snippettheorem}{gauss-lemma-theorem}{Gauss's lemma}
    Let \(A\) be a \uniquefactorizationdomain and let
    \(K=\quotientfield(A)\). Then \(A[x]\) is a
    \uniquefactorizationdomain.

    Moreover, an element \(h\in A[x]\) is prime in \(A[x]\)
    \ifandonlyif exactly one of the following descriptions applies:
    \begin{enumerate}
        \item \(h\in A\) and \(h\) is prime in \(A\);
        \item \(h\) is primitive, \(\polynomialdeg h\geq1\), and \(h\)
        is prime in \(K[x]\).
    \end{enumerate}
\end{snippettheorem}

\begin{snippetproof}{gauss-lemma-theorem-proof}{gauss-lemma-theorem}{Gauss's lemma}
    We first record a normalization fact. For every nonzero
    \(f\in K[x]\), there exists \(\alpha\in K^\times\) such that
    \(\alpha f\in A[x]\) is primitive. Indeed, first multiply by a common
    denominator of the coefficients. The resulting polynomial lies in
    \(A[x]\); dividing it by the common prime powers of its coefficients
    produces a primitive scalar multiple of \(f\).

    We now characterize the prime elements of \(A[x]\).

    Suppose first that \(h\in A\) is prime in \(A\). Coefficientwise
    reduction induces a surjective \ringhomomorphism
    \[
        \varphi\colon A[x]\fromto(A/(h))[x]
    \]
    with kernel \((h)\). Hence
    \[
        A[x]/(h)\ringisomorphic(A/(h))[x].
    \]
    The right-hand side is an \integraldomain because \((h)\) is a prime
    \ideal of \(A\). Thus \((h)\) is a prime ideal of \(A[x]\), so \(h\)
    is prime in \(A[x]\).

    Next, suppose that \(h\in A[x]\) is primitive, has positive degree,
    and is prime in \(K[x]\). If \(h\divides f_1f_2\) in \(A[x]\), then
    the same divisibility holds in \(K[x]\), so \(h\divides f_1\) or
    \(h\divides f_2\) there. Say \(f_1=hg\) for some \(g\in K[x]\). For
    every \primeelement \(p\in A\),
    \[
        \nu_p(g)
        =
        \nu_p(f_1)-\nu_p(h)
        =
        \nu_p(f_1)\geq0.
    \]
    Thus every coefficient of \(g\) belongs to \(A\), so \(g\in A[x]\).
    Therefore \(h\divides f_1\) already in \(A[x]\), and \(h\) is prime
    in \(A[x]\).

    Conversely, let \(h\in A[x]\) be prime. If
    \(\polynomialdeg h=0\), then \(h\in A\), and applying the prime
    divisibility property to constant polynomials shows that \(h\) is
    prime in \(A\).

    Suppose instead that \(\polynomialdeg h\geq1\). A prime element is
    irreducible. If a prime \(p\in A\) divided every coefficient of
    \(h\), then \(h=ph_1\) would be a factorization into two nonunits of
    \(A[x]\), a contradiction. Hence \(h\) is primitive.

    To prove that \(h\) is prime in \(K[x]\), suppose
    \(h\divides FG\) for \(F,G\in K[x]\). Choose
    \(\alpha,\beta\in K^\times\) such that
    \(\alpha F,\beta G\in A[x]\) are primitive. Since \(h\) and
    \((\alpha F)(\beta G)\) are primitive, their quotient in \(K[x]\)
    has nonnegative valuation at every prime of \(A\), and therefore
    belongs to \(A[x]\). Thus
    \[
        h\divides(\alpha F)(\beta G)
        \qquad\text{in }A[x].
    \]
    Primality in \(A[x]\) gives
    \(h\divides\alpha F\) or \(h\divides\beta G\), and hence
    \(h\divides F\) or \(h\divides G\) in \(K[x]\).

    It remains to prove that \(A[x]\) is a
    \uniquefactorizationdomain. Let \(f\in A[x]\) be a nonzero nonunit.
    Since \(K[x]\) is a \principalidealdomain, it is a
    \uniquefactorizationdomain, so
    \[
        f=c\widetilde h_1\cdots\widetilde h_r
    \]
    in \(K[x]\), where \(c\in K^\times\) and the factors are irreducible
    and hence prime. The list of nonconstant factors may be empty.
    Normalize every factor by choosing \(\alpha_i\in K^\times\) so that
    \[
        h_i=\alpha_i\widetilde h_i\in A[x]
    \]
    is primitive. The characterization above shows that each \(h_i\) is
    prime in \(A[x]\). We may write
    \[
        f=a h_1\cdots h_r
    \]
    for some \(a\in K^\times\). For every prime \(p\in A\),
    \[
        \nu_p(a)=\nu_p(f)\geq0,
    \]
    because all the \(h_i\) are primitive. Hence \(a\in A\).
    Factorizing \(a\) into prime elements of \(A\), which remain prime in
    \(A[x]\), gives a factorization of \(f\) into prime elements of
    \(A[x]\). Such a factorization is unique up to order and associates,
    so \(A[x]\) is a \uniquefactorizationdomain.
\end{snippetproof}

\begin{snippetcorollary}{primitive-polynomial-fraction-field-irreducibility}{Irreducibility over the fraction field}
    Let \(A\) be a \uniquefactorizationdomain, let
    \(K=\quotientfield(A)\), and let \(f\in A[x]\) be primitive. Then
    \[
        f\text{ is irreducible in }A[x]
        \ifandonlyif
        f\text{ is irreducible in }K[x].
    \]
\end{snippetcorollary}

\begin{snippetproof}{primitive-polynomial-fraction-field-irreducibility-proof}{primitive-polynomial-fraction-field-irreducibility}{Irreducibility over the fraction field}
    If \(f\) is irreducible in \(K[x]\), then it is prime there because
    \(K[x]\) is a \principalidealdomain. Gauss's lemma makes \(f\) prime,
    and hence irreducible, in \(A[x]\).

    Conversely, suppose that \(f\) is irreducible in \(A[x]\) and that
    \(f=g_1g_2\) in \(K[x]\). Normalize the factors to obtain primitive
    \(h_1,h_2\in A[x]\) and a scalar \(c\in K^\times\) such that
    \[
        f=c h_1h_2.
    \]
    Since \(f,h_1,h_2\) are primitive,
    \(\nu_p(c)=0\) for every prime \(p\in A\), so \(c\in A^\times\).
    This is a factorization in \(A[x]\); irreducibility of \(f\) makes
    one of the \(h_i\) a unit in \(A[x]\), and the corresponding \(g_i\)
    is a unit in \(K[x]\). Thus \(f\) is irreducible in \(K[x]\).
\end{snippetproof}

\end{document}
